\scanpage{56}
\mspage{47}
% archive sheet number: 56

pour toute $A\in\unsure{\Phi|P}_0$.

L'application \unsure{(1)|(1 bis)} peut être vue comme un ``passage au quotient'', obtenu en posant dans $\mathrm{Fig}(\unsure{L|\mathscr{L}})$ formellement
\[
\text{``}X=\varnothing\text{''}
\]
pour tte multistrate $X$ telle que
\[
X^\circ\cap L_0=\varnothing\qquad\text{i.e.}\quad F_{L_0}\mathrel{|\circ|}X
\]
\deleted{Il y a \unsure{danger} \unsure{toutefois}, \unsure{une}}
(mais attention, on conserve les sous-strates $Y$ pour lesquelles $Y^\circ\cap L_0\ne\varnothing$).

\deleted{Il y} On devine que de façon générale, un tel passage au quotient aura un sens (quand on se donne, dans un magasin $M$, une partie qui soit $|\circ|$-fermée, ou un ``$|\circ|$''-support, dont les éléments seront \unsure{égalés} à $0$. Si on néglige $\deleted{\unsure{R}}X$, on ne néglige pas pour autant toutes ses strates, mais du moins tous les $Y$ tels que $Y\mathrel{\overset{\circ}{\ll}}X$ (\unsure{moralement}, tels que $Y^\circ\subset X^\circ$)
% slanted
\margin{NB}

\deleted{\unsure{L'} L'application}

L'application (1 ter) est un isomorphisme, inverse de
\diagram{\begin{tikzcd}[ampersand replacement=\&,column sep=small]
\mathrm{Figures}_{\deleted{\unsure{L_0}}}(L_0)\arrow[r,"\sim"] \& \mathrm{Figures}(L)_{L_0}\arrow[r,hook] \& \mathrm{Figures}(L)
\end{tikzcd}}
et le composé
\[
i_*:\mathrm{Figures}(L_0)\unsure{\relbar\joinrel\relbar|\longrightarrow}\mathrm{Figures}(L)
\]
est un ``plongement'', d'image $\mathrm{Figures}(L)_{L_0}$.
